% Section: P5 MIN-REPORT v2.2 value-correctness mutation stress test (iter 121)
% Closes the AUDITOR-ROBUSTNESS gap by testing the iter-121 value-correctness
% auditor against 8 controlled mutations applied to the n=98 live mega
% corpus. Iter-97/105/113/117 audited manifest CONTENTS at schema/value/content/
% structural layers. None of them stress-tested the AUDITORS THEMSELVES.
\subsection{MIN-REPORT v2.2 value-correctness mutation stress test}
\label{sec:p5-iter121-value-correctness}

A presence-only audit (``is field $X$ present in the manifest?'') can pass
on a manifest whose values are silently wrong. The value-correctness
auditor (10 checks, extending the iter-113 emit-gap audit in
Section~\ref{sec:p5-iter113-v22-recovery}) must be tested against
controlled perturbations to quantify: (a) detection rate for each
mutation class, (b) blind spots (mutations that pass the auditor
unchanged), and (c) the check-fanout structure (which check fires on
which mutation).

\subsubsection{Method}
We load the $n{=}98$ live mega manifests and apply $8$ controlled
mutations per cell. Each mutation corrupts a known MIN-REPORT v2.2 item
by flipping the value to a known-wrong one (preserving format but
breaking semantics): \textbf{M1} cell\_id hash swap; \textbf{M2}
model\_family swap; \textbf{M3} task\_slice swap;
\textbf{M4} $G$ swap; \textbf{M5} temperature swap;
\textbf{M6} seed swap; \textbf{M7} heldout\_split JSON swap;
\textbf{M8} per\_step\_zvf\_path break. We then re-run the
$10$-check value-correctness audit; a mutation is ``caught'' if at
least one check flips from \textsc{pass} to \textsc{fail}.

\subsubsection{Hypotheses}
\begin{itemize}
  \item \textbf{H1 (falsifiable)}: every mutation is caught on $\geq 1$
        check (auditor is complete; no silent corruption).
  \item \textbf{H2 (falsifiable)}: every mutation has a stable dedicated
        check with detection-rate $\geq 0.95$ (one-to-one
        M$_i \leftrightarrow$ C$_j$ mapping).
  \item \textbf{H3 (falsifiable)}: cell\_id mutations flip C01; JSON-body
        mutations flip C07/C10.
  \item \textbf{H4 (falsifiable)}: there exists at least one mutation
        class with detection-rate $0.0$ (auditor has blind spots).
\end{itemize}

\subsubsection{Measured results}

\begin{table}[ht]
\centering
\small
\begin{tabular}{l l r r r l r}
\toprule
mut & name & n\_cells & n\_caught & det\_rate & top\_check & top\_count \\
\midrule
M1 & cell\_id\_swap\_hash        & 98 &  \textbf{0} & \textbf{0.000} & C01 & 0  \\
M2 & model\_family\_swap         & 98 &  \textbf{0} & \textbf{0.000} & C01 & 0  \\
M3 & task\_slice\_swap           & 98 & 98 & \textbf{1.000} & C03 & 98 \\
M4 & G\_swap                     & 98 & 98 & \textbf{1.000} & C04 & 98 \\
M5 & temperature\_swap           & 98 & 98 & \textbf{1.000} & C05 & 98 \\
M6 & seed\_swap                  & 98 & 98 & \textbf{1.000} & C06 & 98 \\
M7 & heldout\_split\_swap        & 98 & 98 & \textbf{1.000} & C07 & 98 \\
M8 & per\_step\_zvf\_path\_break & 98 & 98 & \textbf{1.000} & C10 & 98 \\
\bottomrule
\end{tabular}
\caption{P5 iter-121 mutation stress test (n=98 cells, 8 mutations = 784
evaluations). H1 REFUTED: M1 and M2 are caught on 0/98 cells.}
\label{tab:p5-iter121-mutation}
\end{table}

Per-check fanout across all 784 evaluations: C03, C04, C05, C06, C07, C10
each fire on $98$ mutation-evaluations ($12.5\%$ of the $784$). C01, C02,
C08, C09 fire on $0$ evaluations --- each is the dedicated check for a
mutation class that the current auditor does not catch.

\subsubsection{Verdicts}

\begin{itemize}
  \item \textbf{H1 REFUTED.} M1 and M2 are caught on 0/98 cells. The
        auditor has blind spots. The unmutated corpus passes $100\%$
        on all $10$ checks (per the iter-113 emit-gap baseline),
        but the auditor does not validate the cell\_id hash suffix
        against a known-registry (M1) and does not cross-validate
        filename model\_family against the cells.tsv model\_family
        ground-truth (M2).
  \item \textbf{H2 PARTIALLY PASS.} $6/8$ mutations have a stable
        dedicated check with detection-rate $1.0$ (M3$\to$C03,
        M4$\to$C04, M5$\to$C05, M6$\to$C06, M7$\to$C07,
        M8$\to$C10). The remaining $2$ (M1, M2) have NO check.
  \item \textbf{H3 REFUTED for M1.} M1 (cell\_id hash swap) does NOT
        flip C01 because the cell\_id JSON-basename match passes
        regardless of the hash suffix value. M2 (model\_family swap)
        also does NOT flip C02 because C02 only validates the model
        token against a $2$-element canonical set, not against
        cells.tsv ground-truth.
  \item \textbf{H4 CONFIRMED.} M1 and M2 are silent-corruption blind
        spots with detection-rate $0.0$ across $98$ cells (combined
        $0/196 = 0.0\%$).
\end{itemize}

\subsubsection{Cross-paper coupling}

The value-correctness audit (passes $100\%$ on the unmutated corpus)
and the mutation stress test (REFUTES H1, CONFIRMS H4) jointly
establish: (i) the unmutated corpus is internally consistent
($10/10$ checks pass on $98/98$ cells, strict-pass rate $100.0\%$);
(ii) the auditor has $2$ specific silent-corruption blind spots
(cell\_id hash suffix, model\_family cross-validation); (iii) the
detection fanout is mostly one-to-one (M$_i \leftrightarrow$ C$_j$),
which is a positive structural property for future extension.
\textbf{Operational recommendation}: add \texttt{C11} (cell\_id hash
matches a known-registry hash) and \texttt{C12} (model\_family
filename token matches cells.tsv model\_family) to close the H4
blind spots; wire into a CI gate in \texttt{registry\_validate.py}.

\subsubsection{Recommended next-iter veins}

\textbf{(a)} Compute the auditor's \emph{type-I} false-positive rate
on a corpus with random benign perturbations (single-character typos
in non-key fields). \textbf{(b)} Compute the auditor's \emph{type-II}
miss rate on a corpus with all $8$ mutations applied simultaneously
(does the auditor still find at least one?). \textbf{(c)} Extend the
mutation set to $16$ by adding the $8$ iter-117 ``structural'' blind
spots (e.g., flip cell\_id filename to a non-canonical task\_slice that
breaks \texttt{parse\_filename} entirely).